% ══════════════════════════════════════════════════════════════════════════════
%  EXPERIMENT 3 — DENSITY BY DISPLACEMENT                              (Lesson 9)
%  Content from the A4 master's Lesson 15.
% ══════════════════════════════════════════════════════════════════════════════
\experimentdiv{Density by Displacement}{Volume of an irregular object \textbullet{} Measuring mass and volume \textbullet{} Calculating density \textbullet{} Identifying the material}{9}\label{sec:3}
\lessoncodes{VC2S8U05 \textbullet{} VC2S8I01, I03, I04, I06, I08}

\vocabhead
\begin{ctab}
\begin{tabular}{|SW{4.2cm}|W{11.6cm}|}
  \hline\hd & \thead{Word} & \thead{Definition (copy from the slide)} \\ \hline
  & \scaf{\textbf{Displacement}} & \scaf{Water being pushed out of the way when an object is put into it.} \\ \hline
  & \textbf{Irregular object} & \\ \hline
  & \textbf{Density} & \\ \hline
  & \textbf{Accuracy} & \\ \hline
  & \textbf{Conclusion} & \\ \hline
\end{tabular}
\end{ctab}

\begin{lipanel}
\inquiryli{Measure volume by displacement and mass on a balance; process the data to calculate density.}
\understandli{Use density to identify a material and predict whether it floats.}
\sctitle
\begin{itemize}[leftmargin=1.4em, itemsep=0pt, topsep=1pt]
  \item I can find the volume of an irregular object by displacement.
  \item I can calculate density with $D = m \div V$ and the correct unit.
  \item I can use density to identify a material, and write a conclusion that answers the aim.
\end{itemize}
\end{lipanel}

\begin{skillsfocus}
\focus{QP1}{Predict the densest object.}
\focus{PC4}{Read the cylinder at eye level.}
\focus{PC7}{Lower each object on a thread.}
\focus{PA1}{Record everything in one table.}
\focus{PA3}{Calculate volume, then density.}
\focus{EV2}{Compare with known densities.}
\focus{EV3}{Is an error too big or too small?}
\focus{CO1}{Write a third-person conclusion.}
\end{skillsfocus}

% ─── TECHNIQUE ────────────────────────────────────────────────────────────────
\exppart{Technique spotlight}{Blue}

\begin{technique}{Volume by displacement}{PC7}
A ruler works for a \textbf{regular} object: multiply its length, width and height. An \textbf{irregular} object, like a rock or a key, needs \textbf{displacement}. Lower it into water and the water level rises by exactly the object's volume, because the object and the water cannot be in the same place at the same time.
\begin{enumerate}[leftmargin=1.6em, itemsep=1pt, topsep=3pt]
  \item Half-fill a measuring cylinder and read the volume at eye level.
  \item Tie a thread to the object and lower it in gently until it is completely under water.
  \item Read the new volume. Volume of object = volume after $-$ volume before.
  \item 1\,mL = 1\,cm\textsuperscript{3}, so the volume in mL is the volume in cm\textsuperscript{3}.
\end{enumerate}
\diagDisplacement{50}{62}
\expfigcap{Volume of the rock $=$ 62\,mL $-$ 50\,mL $=$ 12\,mL $=$ 12\,cm\textsuperscript{3}.}
\end{technique}

\begin{technique}{Calculating density and identifying a material}{PA3}
\textbf{Density} $=$ \textbf{mass} $\div$ \textbf{volume}, \quad $D = \dfrac{m}{V}$, \quad in g/cm\textsuperscript{3}. Compare your answer with known densities and choose the closest.

\smallskip
\begin{center}\small
\begin{tabular}{|l|*{8}{c|}}
  \hline\thead{Material} & wood & plastic & glass & aluminium & iron & brass & copper & lead \\ \hline
  \thead{Density (g/cm\textsuperscript{3})} & 0.7 & 1.2 & 2.5 & 2.7 & 7.9 & 8.5 & 8.9 & 11.3 \\ \hline
\end{tabular}
\end{center}
\end{technique}

\begin{workedex}{Density by displacement}
A metal sinker has a mass of 88.5\,g. In a measuring cylinder, the water level rises from 40\,mL to 50\,mL when it is added. Find its density, and identify the metal: aluminium 2.7, iron 7.9 or copper 8.9\,g/cm\textsuperscript{3}.

\medskip
\textbf{Volume:} $50 - 40 = 10\ \text{mL} = 10\ \text{cm}^3$

\medskip
\textbf{Density:} $D = \dfrac{m}{V} = \dfrac{88.5}{10} = 8.85\ \text{g/cm}^3$ \quad --- closest to \textbf{copper}.
\end{workedex}

\begin{wedo}{Density by displacement}
A pebble has a mass of 37.5\,g. The water level rises from 60\,mL to 75\,mL when it is added. Find its density, and decide whether it is more likely to be glass (2.5\,g/cm\textsuperscript{3}) or iron (7.9\,g/cm\textsuperscript{3}).
\workbox{3.5cm}
\end{wedo}

\begin{youdobox}[Density by displacement]
\taskpair{A key has a mass of 34\,g and raises the water from 50\,mL to 54\,mL. Find its density.}{3.5cm}
         {A marble has a mass of 12.5\,g and raises the water from 30\,mL to 35\,mL. Find its density.}{3.5cm}
\end{youdobox}

% ─── PLAN AND SAFETY / METHOD ────────────────────────────────────────────────
\exppart{Plan, safety and method}{Teal}

\textbf{Aim:} To measure the density of several objects and use it to identify what they are made of.

\needspace{10\baselineskip}
\task Look at your objects and feel how heavy each one is for its size. Which do you think is the densest? Explain your prediction. \skr{QP1}
\answerbox{5cm}

\subsubsection*{Equipment}
\begin{multicols}{2}
\begin{itemize}[leftmargin=1.2em, itemsep=0pt, topsep=0pt]
  \item electronic balance
  \item 100\,mL measuring cylinder
  \item ruler
  \item water, paper towel and thread
  \item a rectangular block
  \item four small irregular objects that sink (for example a rock, a metal sinker, a marble and a key)
\end{itemize}
\end{multicols}

\begin{infobox}{Safety}
Lower objects in gently on a thread --- a dropped object can crack the bottom of a glass measuring cylinder. Wipe up spills, and keep water away from the balance.
\end{infobox}

\subsubsection*{Method}
\begin{enumerate}[itemsep=2pt]
  \item Measure the mass of each object on the balance and record it.
  \item For the rectangular block, measure its length, width and height with the ruler, and multiply them to find its volume.
  \item Half-fill the measuring cylinder with water. Read the volume at eye level and record it.
  \item Tie thread to an irregular object and lower it gently into the cylinder until it is completely under water. Do not let it drop.
  \item Read the new volume and record it. Subtract to find the object's volume.
  \item Dry the object, and repeat steps 3--5 for each irregular object.
\end{enumerate}

% ─── RESULTS ─────────────────────────────────────────────────────────────────
\exppart{Results}{Amber}

\subsubsection*{Mass, volume and density of each object \skr{PA1}}
\begin{ctab}
\begin{tabular}{|SW{2.6cm}|Z{1.8cm}|Z{2.2cm}|Z{2.2cm}|Z{2.0cm}|Z{2.4cm}|}
  \hline\hd & \thead{Object} & \thead{Mass (g)} & \thead{Water before (mL)} & \thead{Water after (mL)} & \thead{Volume (cm\textsuperscript{3})} & \thead{Density (g/cm\textsuperscript{3})} \\ \hline
  & \scaf{e.g.\ sinker} & \scaf{88.5} & \scaf{40} & \scaf{50} & \scaf{10} & \scaf{8.85} \\ \hline
  & Block & & \multicolumn{2}{c|}{\small measured with a ruler} & & \\ \hline
  & & & & & & \\ \hline
  & & & & & & \\ \hline
  & & & & & & \\ \hline
  & & & & & & \\ \hline
\end{tabular}
\end{ctab}

\needspace{10\baselineskip}
\task Show your working for the volume and density of your block. \skr{PA3}
\workbox{4cm}

% ─── ANALYSIS ────────────────────────────────────────────────────────────────
\exppart{Analysis and evaluation}{Coral}

\task Identify the material of two of your objects. Use these known densities, and give a reason for each choice: wood 0.7, plastic 1.2, glass 2.5, aluminium 2.7, iron 7.9, brass 8.5, copper 8.9, lead 11.3\,g/cm\textsuperscript{3}. \skr{EV2}
\answerbox{5cm}

\needspace{10\baselineskip}
\task Was your prediction about the densest object correct? Explain using your results. \skr{EV1}
\answerbox{5cm}

\needspace{10\baselineskip}
\task Which of your objects would float in water (density 1.0\,g/cm\textsuperscript{3})? Explain how you can tell from your results.
\answerbox{5cm}

\needspace{14\baselineskip}
\task Evaluate the method. \skr{EV3}
\pt{a}An air bubble stuck to an object as it went into the water. Would this make the density you calculated too big or too small? Explain.
\answerbox{5cm}
\pt{b}Suggest one way to make the volume measurements more accurate.
\answerbox{5cm}

\needspace{10\baselineskip}
\task Write a conclusion for this experiment in the third person. Answer the aim and use data from your table. \skr{CO1}

\scaf{The densities of the objects were measured by displacement. The densest object was \dots}
\answerbox{5.5cm}

\begin{yr8box}
\begin{keepext}
\exttask A cork floats, so you cannot sink it to measure its volume by displacement. Describe a way to get round this, and explain how you would work out the cork's volume.
\answerbox{5cm}
\end{keepext}
\begin{keepext}
\exttask A student uses a 100\,mL measuring cylinder with 10\,mL markings to find the volume of a 2\,cm\textsuperscript{3} button. Explain why her density result is likely to be very inaccurate, and suggest a better approach.
\answerbox{5cm}
\end{keepext}
\begin{keepext}
\exttask An overflow can is filled to its spout. A stone is lowered in and 18\,mL of water flows out into a measuring cylinder. The stone has a mass of 45\,g.
\pt{a}Calculate its density.
\workbox{3cm}
\pt{b}Explain why an overflow can is useful for objects too big for a measuring cylinder.
\answerbox{5cm}
\end{keepext}
\end{yr8box}

% ─── SKILL CHECK ─────────────────────────────────────────────────────────────
\exppart{Skill check}{NavyBlue}

\begin{skillcheck}
\scrow{QP1}{I can make and explain a prediction.}
\scrow{PC4}{I read the measuring cylinder at eye level before and after each object.}
\scrow{PC7}{I lowered each object on a thread and used the equipment safely.}
\scrow{PA1}{I recorded mass, volumes and density in one table with units.}
\scrow{PA3}{I can calculate volume by difference and density with $D = m \div V$.}
\scrow{EV2}{I can compare my densities with known values to identify a material.}
\scrow{EV3}{I can say whether an error makes my result too big or too small.}
\scrow{CO1}{I can write a conclusion in the third person that answers the aim.}
\end{skillcheck}

\task Fill in the \emph{End of unit} column of your skills tracker at the front of this booklet. Which skill has improved most since the Viscosity Race, and which one will you work on next?
\answerbox{5cm}

\end{document}
